% Источники: 03_Карточки_доменов/Сессии_задач.md
% Генератор: docs/build_pdf.py
% Шаблон печатного представления. Содержимое подставляет docs/build_pdf.py.
\documentclass[12pt,a4paper]{article}
\usepackage{iftex}
\ifPDFTeX
  \usepackage[utf8]{inputenc}
  \usepackage[T2A]{fontenc}
  \usepackage{newtxtext,newtxmath}
\else
  \usepackage{fontspec}
  \IfFontExistsTF{Times New Roman}
    {\setmainfont{Times New Roman}}
    {\setmainfont{Libertinus Serif}}
  \IfFontExistsTF{Arial}{\setsansfont{Arial}}{\setsansfont{Latin Modern Sans}}
  \IfFontExistsTF{Menlo}{\setmonofont{Menlo}[Scale=MatchLowercase]}
    {\setmonofont{Latin Modern Mono}[Scale=MatchLowercase]}
\fi
\usepackage[russian]{babel}
\usepackage{geometry}
\geometry{left=25mm,right=15mm,top=20mm,bottom=20mm}
\usepackage{setspace}
\onehalfspacing
\usepackage{indentfirst}
\setlength{\parindent}{1.25cm}
\setlength{\parskip}{2pt}
\setlength{\emergencystretch}{3em}
\usepackage{graphicx,booktabs,longtable,array,tabularx,calc,etoolbox}
\usepackage{enumitem}
\setlist[itemize]{leftmargin=1.25cm,itemsep=2pt,topsep=4pt}
\setlist[enumerate]{leftmargin=1.25cm,itemsep=2pt,topsep=4pt}
\newcommand{\tightlist}{\setlength{\itemsep}{0pt}\setlength{\parskip}{0pt}}
\AtBeginEnvironment{longtable}{\small\singlespacing}
\setlength{\LTpre}{8pt}
\setlength{\LTpost}{8pt}
\setlength{\tabcolsep}{5pt}
\setlength{\LTleft}{0pt}
\setlength{\LTright}{0pt}
\usepackage{xurl}
\usepackage[unicode,colorlinks=true,linkcolor=black,urlcolor=blue]{hyperref}
\usepackage{bookmark}
\hypersetup{pdftitle={Карточка домена: сессии задач},pdfauthor={Анохин Г.А.},pdfsubject={Криптотренажёр mephi-12}}
\urlstyle{same}
\pagestyle{plain}
\setcounter{secnumdepth}{3}
\setcounter{tocdepth}{2}

\begin{document}
\hypersetup{pageanchor=false}
\begin{titlepage}
  \thispagestyle{empty}
  \centering
  {\large НИЯУ МИФИ\par}
  \vspace*{3.5cm}
  {\LARGE\bfseries Карточка домена: сессии задач\par}
  \vspace{0.5cm}
  {\large Карточка домена\par}
  \vspace{0.6cm}
  {\large Платформа для изучения и практики курса по криптографии\par}
  {\large Криптотренажёр mephi-12\par}
  \vspace{0.7cm}
  {\large Версия 0.3\par}
  {\large Редакция от 07.11.2026\par}
  \vspace{1cm}
  \begin{tabularx}{\textwidth}{@{}p{3.4cm}Xp{4cm}@{}}
    \toprule
    Роль & ФИО & Подпись / Дата \\
    \midrule
    Разработал & Анохин Г.А. & \\
    \addlinespace[6pt]
    Согласовано & & \\
    \addlinespace[6pt]
    Утверждено & & \\
    \bottomrule
  \end{tabularx}
  \vfill
  {\large Москва, 2026\par}
\end{titlepage}
\hypersetup{pageanchor=true}

\section*{Лист регистрации изменений}
\addcontentsline{toc}{section}{Лист регистрации изменений}
\begingroup\small\singlespacing
\noindent
\begin{tabularx}{\textwidth}{@{}p{1.2cm}p{2cm}p{2.5cm}Xp{1.8cm}@{}}
\toprule
Версия & Дата & Автор & Содержание изменения & Утвердил \\
\midrule
0.3 & 07.11.2026 & Анохин Г.А. & Уточнены правила и контракты; расширено описание сценариев; подготовлена печатная редакция. & \\
\bottomrule
\end{tabularx}
\endgroup
\clearpage
\begingroup\small\singlespacing
\tableofcontents
\endgroup
\clearpage

\hypertarget{карточка-домена-сессии-задач}{}

Версия: 0.3. Дата: 07.11.2026. Владелец: Анохин Г.А. Документ определяет поведение домена сессий задач.

\section{Назначение и границы}\label{ux43dux430ux437ux43dux430ux447ux435ux43dux438ux435-ux438-ux433ux440ux430ux43dux438ux446ux44b}

\hypertarget{назначение-и-границы}{}

Домен создаёт пользователю набор задач по доступной конфигурации, принимает ответы, определяет результат, предоставляет его для чтения и контролирует срок. Для внешних участников он описывается через операции создания сессии и отправки ответа. Устройство генераторов и алгоритмы проверки относятся к криптографическим доменам.

{\def\LTcaptype{none} % do not increment counter
\begin{longtable}[]{@{}
  >{\raggedright\arraybackslash}p{(\linewidth - 2\tabcolsep) * \real{0.3200}}
  >{\raggedright\arraybackslash}p{(\linewidth - 2\tabcolsep) * \real{0.6800}}@{}}
\toprule\noalign{}
\begin{minipage}[b]{\linewidth}\raggedright
Смежная область
\end{minipage} & \begin{minipage}[b]{\linewidth}\raggedright
Что требуется домену
\end{minipage} \\
\midrule\noalign{}
\endhead
\bottomrule\noalign{}
\endlastfoot
Аутентификация & Подтверждённая личность пользователя и полномочия \\
Группы и конфигурации & Группа пользователя, включённые конфигурации и состав выбранного набора \\
Криптографические домены & Генерация условия и предметная проверка для каждого поддерживаемого типа \\
Хранение & Состояния сессии и задач, принадлежность и серверное время создания \\
\end{longtable}
}

\section{Сценарии}\label{ux441ux446ux435ux43dux430ux440ux438ux438}

\hypertarget{сценарии}{}

При создании набора система проверяет доступность конфигурации для группы и отсутствие другой незавершённой сессии пользователя. Затем формирует задачи и возвращает созданную сессию. Её состав определяется выбранной конфигурацией.

При отправке ответа система устанавливает пользователя и задачу, проверяет принадлежность, состояние и своевременность запроса. После проверки формата решение передаётся проверяющему соответствующего типа. Вердикт определяет итог задачи; когда завершены все задачи, определяется итог сессии.

\section{Срок и попытка}\label{ux441ux440ux43eux43a-ux438-ux43fux43eux43fux44bux442ux43aux430}

\hypertarget{срок-и-попытка}{}

Обозначения: \texttt{t0} --- серверное время создания сессии; \texttt{D\ =\ t0\ +\ 20\ минут} --- общий срок; \texttt{t\_in} --- время получения сервером полного запроса с ответом. Все задачи сессии используют один \texttt{D}.

{\def\LTcaptype{none} % do not increment counter
\begin{longtable}[]{@{}
  >{\raggedright\arraybackslash}p{(\linewidth - 2\tabcolsep) * \real{0.3200}}
  >{\raggedright\arraybackslash}p{(\linewidth - 2\tabcolsep) * \real{0.6800}}@{}}
\toprule\noalign{}
\begin{minipage}[b]{\linewidth}\raggedright
Условие
\end{minipage} & \begin{minipage}[b]{\linewidth}\raggedright
Правило
\end{minipage} \\
\midrule\noalign{}
\endhead
\bottomrule\noalign{}
\endlastfoot
\texttt{t\_in\ \textless{}\ D} & Ответ своевременный. При выполнении остальных условий он допускается к проверке \\
\texttt{t\_in\ \textgreater{}=\ D} & Новый ответ отклоняется по сроку \\
Своевременная проверка закончилась после \texttt{D}, но до \texttt{t\_in\ +\ 60\ секунд} & Её вердикт принимается; истечение общего срока во время работы сервера не отменяет своевременный ответ \\
Получен вердикт «верно» или «неверно» & Расходуется единственная попытка задачи; устанавливается итоговое состояние \\
Ошибка формата или нет обязательного поля & Предметная попытка не расходуется; исправленный запрос допустим только при его получении до \texttt{D} \\
Технический сбой без вердикта & Предметная попытка не расходуется; новый запрос допустим только до \texttt{D} \\
После \texttt{D} нет результата и нет своевременного ответа, ожидающего проверки & Незавершённая задача получает \texttt{FAILED} по просрочке \\
\end{longtable}
}

Предметная попытка завершается фиксацией вердикта в хранилище. Обновление страницы, ошибки и повторные обращения не изменяют \texttt{t0} и \texttt{D}. Если своевременная проверка завершилась техническим сбоем уже после \texttt{D}, повторный ответ не допускается; задача завершается по просрочке. Основание завершения --- общий срок; предметная попытка остаётся неиспользованной.

Пока своевременный ответ обрабатывается, задаче не назначается просрочка. На обработку ответа отводится 60 секунд от \texttt{t\_in}, включая ожидание и восстановление после перезапуска. По окончании этого времени незавершённая обработка прекращается с техническим отказом без расхода предметной попытки. Новый ответ допустим только до \texttt{D}; если общий срок уже наступил, задача завершается по просрочке.

Принятый запрос сохраняется до запуска предметного вычисления. После перезапуска обработка продолжается по тем же входным данным, если её срок ещё не истёк. Время поступления и оба срока не пересчитываются. Фиксация вердикта завершает предметную попытку; повтор вычисления после сбоя не должен создавать второй зафиксированный результат.

Одновременные обращения не должны приводить к двум завершённым предметным попыткам одной задачи. Программные обязанности по согласованию проверки и планировщика определены в \href{../07_\%D0\%9F\%D1\%80\%D0\%BE\%D0\%B3\%D1\%80\%D0\%B0\%D0\%BC\%D0\%BC\%D0\%BD\%D1\%8B\%D0\%B5_\%D1\%82\%D1\%80\%D0\%B5\%D0\%B1\%D0\%BE\%D0\%B2\%D0\%B0\%D0\%BD\%D0\%B8\%D1\%8F/DOC-06.2_\%D0\%A1\%D0\%B5\%D1\%81\%D1\%81\%D0\%B8\%D0\%B8_\%D0\%B7\%D0\%B0\%D0\%B4\%D0\%B0\%D1\%87_\%D0\%B8_\%D0\%BF\%D0\%BB\%D0\%B0\%D0\%BD\%D0\%B8\%D1\%80\%D0\%BE\%D0\%B2\%D1\%89\%D0\%B8\%D0\%BA.pdf\#nameddest=doc-062-\%D1\%81\%D0\%B5\%D1\%81\%D1\%81\%D0\%B8\%D0\%B8-\%D0\%B7\%D0\%B0\%D0\%B4\%D0\%B0\%D1\%87-\%D0\%B8-\%D0\%BF\%D0\%BB\%D0\%B0\%D0\%BD\%D0\%B8\%D1\%80\%D0\%BE\%D0\%B2\%D1\%89\%D0\%B8\%D0\%BA}{DOC-06.2}.

\section{Итог сессии}\label{ux438ux442ux43eux433-ux441ux435ux441ux441ux438ux438}

\hypertarget{итог-сессии}{}

Сессия остаётся \texttt{NEW}, пока существует задача без итогового состояния, в том числе своевременный ответ в обработке. После завершения всех задач сессия получает \texttt{SOLVED}, если все задачи решены верно, и \texttt{FAILED} в остальных случаях. Ранее определённые результаты задач сохраняются при завершении по сроку.

\section{Связанные материалы}\label{ux441ux432ux44fux437ux430ux43dux43dux44bux435-ux43cux430ux442ux435ux440ux438ux430ux43bux44b}

\hypertarget{связанные-материалы}{}

В \href{../04_\%D0\%A1\%D0\%B8\%D1\%81\%D1\%82\%D0\%B5\%D0\%BC\%D0\%BD\%D1\%8B\%D0\%B5_\%D1\%82\%D1\%80\%D0\%B5\%D0\%B1\%D0\%BE\%D0\%B2\%D0\%B0\%D0\%BD\%D0\%B8\%D1\%8F/DOC-05_\%D0\%A1\%D0\%B8\%D1\%81\%D1\%82\%D0\%B5\%D0\%BC\%D0\%BD\%D1\%8B\%D0\%B5_\%D1\%82\%D1\%80\%D0\%B5\%D0\%B1\%D0\%BE\%D0\%B2\%D0\%B0\%D0\%BD\%D0\%B8\%D1\%8F.pdf\#nameddest=doc-05-\%D1\%82\%D1\%80\%D0\%B5\%D0\%B1\%D0\%BE\%D0\%B2\%D0\%B0\%D0\%BD\%D0\%B8\%D1\%8F-\%D0\%BA-\%D1\%81\%D0\%B5\%D1\%81\%D1\%81\%D0\%B8\%D1\%8F\%D0\%BC-\%D0\%B7\%D0\%B0\%D0\%B4\%D0\%B0\%D1\%87}{системных требованиях} определены обязательства системы. \href{../07_\%D0\%9F\%D1\%80\%D0\%BE\%D0\%B3\%D1\%80\%D0\%B0\%D0\%BC\%D0\%BC\%D0\%BD\%D1\%8B\%D0\%B5_\%D1\%82\%D1\%80\%D0\%B5\%D0\%B1\%D0\%BE\%D0\%B2\%D0\%B0\%D0\%BD\%D0\%B8\%D1\%8F/DOC-06.2_\%D0\%A1\%D0\%B5\%D1\%81\%D1\%81\%D0\%B8\%D0\%B8_\%D0\%B7\%D0\%B0\%D0\%B4\%D0\%B0\%D1\%87_\%D0\%B8_\%D0\%BF\%D0\%BB\%D0\%B0\%D0\%BD\%D0\%B8\%D1\%80\%D0\%BE\%D0\%B2\%D1\%89\%D0\%B8\%D0\%BA.pdf\#nameddest=doc-062-\%D1\%81\%D0\%B5\%D1\%81\%D1\%81\%D0\%B8\%D0\%B8-\%D0\%B7\%D0\%B0\%D0\%B4\%D0\%B0\%D1\%87-\%D0\%B8-\%D0\%BF\%D0\%BB\%D0\%B0\%D0\%BD\%D0\%B8\%D1\%80\%D0\%BE\%D0\%B2\%D1\%89\%D0\%B8\%D0\%BA}{Программные требования} распределяют их между компонентами. \href{../07_\%D0\%9F\%D1\%80\%D0\%BE\%D0\%B3\%D1\%80\%D0\%B0\%D0\%BC\%D0\%BC\%D0\%BD\%D1\%8B\%D0\%B5_\%D1\%82\%D1\%80\%D0\%B5\%D0\%B1\%D0\%BE\%D0\%B2\%D0\%B0\%D0\%BD\%D0\%B8\%D1\%8F/DOC-06.2_\%D0\%A1\%D0\%B5\%D1\%81\%D1\%81\%D0\%B8\%D0\%B8_\%D0\%B7\%D0\%B0\%D0\%B4\%D0\%B0\%D1\%87_\%D0\%B8_\%D0\%BF\%D0\%BB\%D0\%B0\%D0\%BD\%D0\%B8\%D1\%80\%D0\%BE\%D0\%B2\%D1\%89\%D0\%B8\%D0\%BA.pdf\#nameddest=doc-062-\%D0\%BA\%D0\%BE\%D0\%BD\%D1\%82\%D1\%80\%D0\%B0\%D0\%BA\%D1\%82\%D1\%8B-\%D1\%81\%D0\%B5\%D1\%81\%D1\%81\%D0\%B8\%D0\%B9-\%D0\%B7\%D0\%B0\%D0\%B4\%D0\%B0\%D1\%87}{Контракты} задают параметры операций, форматы и отказы API.

История: 05.11.2026 --- сохранены общий срок и одна предметная попытка, уточнена своевременность ответа. 06.11.2026 --- определены предел обработки и восстановление после перезапуска.


\end{document}
